% Auto-generated person page
\documentclass[10pt, a4paper]{book}
\usepackage[margin=1in]{geometry}
\usepackage{array}
\usepackage{xcolor}
\usepackage{ragged2e}
\usepackage{graphicx}
\newcommand{\birthsymbol}{\ensuremath{\star}}
\newcommand{\deathsymbol}{\ensuremath{\dagger}}

\providecommand{\recordchapter}{2}

\begin{document}
\thispagestyle{empty}
\newpage

\recordchapter{root}{Katarzyna Komasara}

\begin{tabular}{@{}p{0.38\textwidth} p{0.58\textwidth}}
  \fbox{%
    \parbox[c][6cm][c]{\linewidth}{%
      \centering      
        {}
    }%
  } &
  \begin{minipage}[t]{\linewidth}
    \vspace{0pt}
    \Huge \textbf{Katarzyna}\\[0.3em]
    \LARGE \textbf{Komasara}\\[0.3em]
    \small \textit{}\\[0.3em]

\begin{tabular}{@{}l l}
Ojciec: & {Jan Komasara} \\ 
Matka: & {Marianna Bielaska} \\
\end{tabular}
  \end{minipage} \\ 
\end{tabular}

\noindent
\begin{tabular}{@{}l l l}
  \birthsymbol & \textbf{Data Narodzin:} & 1878, Komorów \\ 
  \deathsymbol & \textbf{Data Śmierci: } & 1907.08.15, Gace Słupieckie, Zborówek \\ 

\end{tabular}
\vspace{1.5em}

\begin{tabular}{@{}l l l l l l}
& \textbf{Data} & \textbf{Miejsce} & \textbf{Wiek} & \textbf{Opis} & \textbf{Osoba}\\
\hline
    \birthsymbol & 1878 & Komorów & 0 dni & Narodziny & \\
\deathsymbol & 1890.03.27 & Zborówek & 12 lat & Śmierć (Ojciec) & Jan Komasara\\
\char"26AD & 1894 &  & 16 lat & Ślub & \\
\birthsymbol & 1898 & Szczebrzusz & 20 lat & Narodziny (Córka) & Marianna Stąpór\\
\birthsymbol & 1900 &  & 22 lat & Narodziny (Córka) & Anna Stąpór\\
\birthsymbol & 1902.05.11 & Szczebrzusz & 24 lat & Narodziny (Syn) & Jan Stąpór\\
\birthsymbol & 1905 & Szczebrzusz & 27 lat & Narodziny (Córka) & Agnieszka Stąpór\\
\deathsymbol & 1907.08.15 & Gace Słupieckie, Zborówek & 29 lat & Śmierć & \\
\end{tabular}

\vspace{1.5em}
\newpage
        \begin{center}
        \includegraphics[width=\textwidth]{/Users/kadasie/PRIV_PRIV/GrampsBook/src/pages/pageJW/assets/pictures/1025111f3b6876d2b7b367486cbb.png}
        \newline
        {Akt urodzenia Katarzyna Komasara}
        \newline   \end{center} \newpage
        \begin{center}
        \includegraphics[width=\textwidth]{/Users/kadasie/PRIV_PRIV/GrampsBook/src/pages/pageJW/assets/pictures/101ed9add2d16c9138125a0e6de6.png}
        \newline
        {Akt Małżeństwa Wawrzyniec Stąpór, Katarzyna Komasara}
        \newline 
            \fbox{
            \parbox{\linewidth}{19
Komorów

Działo się we wsi Zborówek, dnia 7/19 listopada 1894 roku, o godzinie 11 w dzień.

Wiadomo czynimy, że w obecności świadków, Bartłomieja Wojtowicz, 50 lat i Sebastiana Fortuny, 70 lat mających, włościan ze wsi Komorowa, na dniu dzisiejszym zawarte zostało religijnie małżeństwo, między:

Wawrzyńcem Stopór, włościaninem, kawalerem, 24 lata mającym, urodzonym i zamieszkałym we wsi Gace, w parafii Szczebrzusz, synem zmarłego Tadeusza i Marianny z domu Skiba

i

Katarzyną Komasara, włościanką, 17 lat mającą, urodzoną i zamieszkałą w miejscowej wsi Komorów, córką Jana i zmarłej Marianny z domu Bielaskiej.

Małżeństwo to poprzedziły 3 zapowiedzi, wygłoszone w kościołach parafialnych, miejscowym zboróweckim i szczebrzuskim, 23, 30 października i 6 listopada bieżącego roku.

Nowo zaślubieni oświadczyli, że nie zawarli żadnej umowy przedślubnej.

Religijna ceremonia zaślubin sprawowana była przez ks. Waleriana Bienieckiego, proboszcza parafii Zborówek.

Akt ten odczytano obecnym, niepiśmiennym, przez Nas tylko podpisany został.

Ks. Walerian Bieniecki.}}
              \end{center} \newpage
        \begin{center}
        \includegraphics[width=\textwidth]{/Users/kadasie/PRIV_PRIV/GrampsBook/src/pages/pageJW/assets/pictures/1024b70818b611d9d82c1a0c52b2.png}
        \newline
        {Akt Zgonu Katarzyna Stąpór}
        \newline 
            \fbox{
            \parbox{\linewidth}{30
Katarzyna Stopór
Gace

Działo się w Szczebrzuszu, dnia 4/17 sierpnia 1907 roku, o godzinie 10 z rana.

Stawili się Marcin Pyrz, 45 lat i Józef Skiba, 46 lat mający, włościanie z Gaców i oświadczyli, że w Gacach dnia 2/15 tegoż miesiąca i roku, o godzinie 10 z rana zmarła Katarzyna Stopór, 29 lat mająca, pozostawiwszy po sobie owdowiałego małżonka, Wawrzyńca Stopór, córka Jana i Marianny z domu Bielaskiej, małżonków Komasara.

Po przekonaniu się naocznie o zejściu Katarzyny Stopór Akt ten odczytano deklarującym, niepiśmiennym i przez Nas tylko podpisany został.

Ks. - podpis mało wyraźny.}}
              \end{center} 


\end{document}
